\documentclass[11pt]{article}
\usepackage[T1]{fontenc}
\usepackage{lmodern,amsmath,amssymb,amsthm,mathtools}
\usepackage[margin=1.05in]{geometry}
\usepackage[hidelinks]{hyperref}
\newtheorem{theorem}{Theorem}[section]
\newtheorem{lemma}[theorem]{Lemma}
\newtheorem{proposition}[theorem]{Proposition}
\newtheorem{corollary}[theorem]{Corollary}
\newcommand{\R}{\mathbb R}\newcommand{\C}{\mathbb C}
\newcommand{\T}{\mathbb T}\newcommand{\D}{\mathbb D}
\newcommand{\Ent}{\operatorname{Ent}}
\newcommand{\code}[1]{{\small\nolinkurl{#1}}}
\title{An explicit nonreversible Cauchy semigroup\protect\\
on the circle: kernel, generator and entropy}
\author{Project Stardust}
\date{Reviewed editorial draft, October 29, 2026}
\begin{document}
\maketitle
\begin{abstract}
A fractional circle generator with sine drift admits an
explicit transition kernel obtained by harmonic extension
along a rational disk flow. We construct its conservative
Feller semigroup and invariant wrapped Cauchy law, identify
the full generator graph on continuous functions and its
Laplace resolvent, and prove graph density of time-smoothed
trigonometric data. For every initial probability law the
positive-time forward density has finite relative entropy.
Its exact entropy production is a singular jump-Bregman
integral, even though the generator is nonsymmetric under
its invariant law. An elementary nonsharp modified
logarithmic-Sobolev bound yields exponential decay from
initial time, with infinite initial entropy allowed.
The proof avoids assigning an unproved backward-generator
domain to a logarithmic density. The construction belongs
to classical harmonic and semigroup theory and is related
to known Cauchy-driven phase dynamics. Exact statements
are mapped to a frozen Lean development; novelty is unresolved.
\end{abstract}

\section{The circle flow and the conclusions}
Let $\T=\{z\in\C:|z|=1\}$ and let $m$ be normalized
Haar measure, so in angle coordinates $dm=d\theta/(2\pi)$.
For $a\in\D$, the Poisson kernel is
\[
 p_a(w)=\frac{1-|a|^2}{|w-a|^2},\qquad w\in\T.
\]
It is a positive probability density with respect to $m$.
Fix a real parameter $|q|<1$, and put
\begin{equation}\label{eq:parameters}
 r=\frac{2q}{1-q^2},\qquad
 \alpha=\frac{1+q^2}{1-q^2},\qquad
 c_q=\frac{1-|q|}{2(1+|q|)}.
\end{equation}
For $0\leq\lambda\leq1$ define
\begin{equation}\label{eq:flow}
 F_\lambda(z)=
 \frac{q(1-\lambda)+(q^2+\lambda)z}
      {1+q^2\lambda+q(1-\lambda)z},\qquad
 \phi_t=F_{e^{-\alpha t}}\quad(t\geq0).
\end{equation}
For $t>0$ define the transition probability
\begin{equation}\label{eq:kernel}
 K_t(z,dw)=p_{\phi_t(z)}(w)\,m(dw),\qquad
 K_0(z,\cdot)=\delta_z,
\end{equation}
and $P_tf(z)=\int f(w)K_t(z,dw)$. Our invariant reference
law is $\nu_q(dw)=p_q(w)m(dw)$.

\begin{theorem}\label{thm:main}
The kernels \eqref{eq:kernel} form a conservative Markov
semigroup. On $C(\T;\C)$, $P_t$ is a strongly continuous
positive contraction semigroup with invariant probability
$\nu_q$. Its closed generator $A$ acts on trigonometric
polynomials as
\begin{equation}\label{eq:generator}
 A=-|D_\theta|-r\sin\theta\,\partial_\theta,
 \qquad |D_\theta|e^{in\theta}=|n|e^{in\theta}.
\end{equation}
Its full graph is characterized by
\begin{equation}\label{eq:graph}
 (f,g)\in\operatorname{graph}A
 \quad\Longleftrightarrow\quad
 P_tf-f=\int_0^tP_sg\,ds\quad\hbox{for every }t\geq0.
\end{equation}
Time averages of trigonometric polynomials are graph
dense. For $\beta>0$ the positive Laplace resolvent
$R_\beta f=\int_0^\infty e^{-\beta t}P_tf\,dt$
is the inverse of $\beta-A$, with norm at most $1/\beta$.

For every Borel probability measure $\mu$, write
$\mu_t=\mu K_t$. At every $t>0$ it has a positive
continuous density $f_t=d\mu_t/d\nu_q$, and
\begin{equation}\label{eq:entropyderivative}
 \frac{d}{dt}D(\mu_t\Vert\nu_q)
 =-\int_{\T^2}\frac{2\,B(f_t(z),f_t(w))}
 {|w-z|^2p_q(w)}\,\nu_q(dz)\nu_q(dw),
\end{equation}
where $B(a,b)=a\log(a/b)-a+b$ and the diagonal integrand
is defined to be zero. The integral is finite, and
\begin{equation}\label{eq:entropybound}
 D(\mu_t\Vert\nu_q)\leq e^{-c_qt}D(\mu\Vert\nu_q),
 \qquad t\geq0,
\end{equation}
in the extended nonnegative reals. The invariant law is
unique. For $q\ne0$ the generator is not Hermitian even
on its finite Fourier domain in $L^2(\nu_q)$.
\end{theorem}

Here $D(\mu\Vert\nu)=\int f\log f\,d\nu$ when
$\mu=f\nu$ and this integral is finite, with the usual
extended-value definition otherwise. The constant $c_q$
is a convenient lower bound, not a sharp entropy constant.
The generator graph is on continuous functions with the
supremum norm. We will not identify it with an unproved
Sobolev domain or replace the smoothed graph core by an
unsmoothed-core assertion.

\section{Disk geometry and harmonic transition kernels}
The denominator in \eqref{eq:flow} has modulus at least
$1+q^2\lambda-|q|(1-\lambda)>0$ on the closed disk.
Writing its numerator and denominator as $N$ and $H$,
direct multiplication gives
\begin{align}\label{eq:defect}
 |H|^2-|N|^2={}&(1-\lambda)^2(1-q^2)|1+qz|^2\notag\\
 &+\lambda^2(1+q^2)^2(1-|z|^2)\notag\\
 &+2\lambda(1-\lambda)(1+q^2)(1-q^2|z|^2).
\end{align}
All terms are nonnegative on $\overline\D$, and the first
is strictly positive for $\lambda<1$. Thus $F_\lambda$
maps the closed disk strictly into $\D$ when $\lambda<1$,
and $F_1(z)=z$. Each map is holomorphic on a neighborhood
of the closed disk.

The matrix defining this fractional-linear map is
\[
 M_\lambda=\begin{pmatrix}q^2+\lambda&q(1-\lambda)\\
 q(1-\lambda)&1+q^2\lambda\end{pmatrix}.
\]
Its eigenvectors $(q,1)$ and $(1,-q)$ have eigenvalues
$1+q^2$ and $\lambda(1+q^2)$, respectively. Hence
$M_\lambda M_\eta=(1+q^2)M_{\lambda\eta}$, proving
\begin{equation}\label{eq:composition}
 \phi_{s+t}=\phi_s\circ\phi_t,\qquad
 \phi_0(z)=z,\qquad \phi_t(q)=q.
\end{equation}
Differentiating \eqref{eq:flow} at time zero gives
\begin{equation}\label{eq:velocity}
 \left.\partial_t\phi_t(z)\right|_{t=0}
   =V(z):=-z+\frac{q}{1-q^2}(1-z^2).
\end{equation}
The denominator lower bound makes this differentiation
uniform on the closed disk. It also gives joint continuity
in $t,z$ and uniform convergence $\phi_t(z)\to z$ as
$t\downarrow0$.

For $f\in C(\T;\C)$, let $u_f$ be its harmonic Poisson
extension to the disk, continuous up to the boundary.
Then
\[
 P_tf(z)=u_f(\phi_t(z)).
\]
The function $u_f\circ\phi_s$ is harmonic in $\D$ and
continuous on its closure. Its boundary values are
$P_sf$. The uniqueness of the Dirichlet problem therefore
gives $u_{P_sf}=u_f\circ\phi_s$, whence
\[
 P_tP_sf=u_f\circ\phi_s\circ\phi_t=P_{s+t}f.
\]
Equality on continuous functions determines finite Borel
measures on the compact circle. This proves the exact
Chapman--Kolmogorov identity for \eqref{eq:kernel}, not
just agreement of a formal generator.

The Poisson kernel is nonnegative with integral one,
so $P_t$ is positive, preserves constants, and is a
supremum-norm contraction. Continuity of $u_f\circ\phi_t$
proves the Feller property. For strong continuity at zero,
use uniform continuity of $u_f$ on $\overline\D$ and
uniform convergence of $\phi_t$ to the identity.
One can also verify boundary continuity of the Poisson
extension directly: split its integral into a small arc
where $f$ oscillates little and its complement, where
$p_a(w)\leq(1-|a|^2)/(\varepsilon-|a-z|)^2$ tends
uniformly to zero as $a\to z\in\T$.

Finally $\int_\T f\,d\nu_q=u_f(q)$, so
\[
 \int_\T P_tf\,d\nu_q=u_{P_tf}(q)
 =u_f(\phi_t(q))=u_f(q).
\]
This proves invariance. The kernels are Borel measurable
in their initial point because their positive-time
densities are jointly continuous, and at zero they are
the measurable identity kernel.

\section{The closed generator and the precise graph core}
For $n\geq0$, the Poisson moment identity gives
\begin{equation}\label{eq:modes}
 P_t(z^n)=\phi_t(z)^n,\qquad
 P_t(\bar z^n)=\overline{\phi_t(z)}^{\,n}.
\end{equation}
Uniform differentiation using \eqref{eq:velocity} shows
that these modes belong to the continuous-function
generator domain. For $n\geq1$,
\[
 Az^n=-nz^n+\frac{nr}{2}(z^{n-1}-z^{n+1}),
\]
with the conjugate formula for negative frequencies and
$A1=0$. Since $\partial_\theta z^n=in z^n$,
these are exactly \eqref{eq:generator}. This statement
concerns actual supremum-norm difference quotients.

Define $f\in D(A)$ when
$h^{-1}(P_hf-f)$ has a supremum-norm limit $g$ as
$h\downarrow0$, and set $Af=g$. The limit is unique.

\begin{proposition}\label{prop:graph}
The equivalence \eqref{eq:graph} holds, and its right-hand
side defines a closed linear graph.
\end{proposition}
\begin{proof}
If $Af=g$, the semigroup identity gives the right derivative
$\partial_t^+P_tf=P_tg$. For $t>0$ the left difference
quotient is
$P_{t-h}(P_hf-f)/h$, also tending to $P_tg$ by
contraction and strong continuity. Thus the orbit is
differentiable for $t>0$ with continuous derivative
$P_tg$, and the vector-valued fundamental theorem of
calculus gives \eqref{eq:graph}.

Conversely, divide that identity by $t$ and let
$t\downarrow0$. Strong continuity gives
$t^{-1}\int_0^tP_sg\,ds\to g$, so $f\in D(A)$ and
$Af=g$.
For a convergent sequence of graph pairs $(f_n,g_n)$,
contraction bounds the difference of the integrals by
$t\|g_n-g\|_\infty$. Pass to the limit in
\eqref{eq:graph}; the limiting pair is again in the graph.
Linearity follows from the same identity.
\end{proof}

For $h>0$ define the time average
\begin{equation}\label{eq:average}
 J_h f=\frac1h\int_0^hP_sf\,ds.
\end{equation}
A shift of the integration interval gives
\begin{equation}\label{eq:averagegraph}
 J_hf\in D(A),\qquad AJ_hf=(P_hf-f)/h.
\end{equation}
Indeed $(P_tJ_hf-J_hf)/t$ is $1/h$ times the
difference of the averages of $P_sf$ on $[h,h+t]$
and $[0,t]$. Letting $t\downarrow0$ proves the formula.
Since $J_hf\to f$, the generator domain is dense.
If $f\in D(A)$, \eqref{eq:graph} also gives
$AJ_hf=J_hAf\to Af$, so these averages approximate in
the full graph norm.

Let $\mathcal P$ be the real-and-complex trigonometric
polynomials. They are uniformly dense in $C(\T;\C)$.
For fixed $h>0$, the map
\[
 f\longmapsto \bigl(J_hf,(P_hf-f)/h\bigr)
\]
is continuous into the product of supremum-norm spaces:
its two component norms are bounded by $1$ and $2/h$.
For a graph pair $(f,Af)$, first choose $h$ small enough
that its averaged pair is close; then choose
$p\in\mathcal P$ close enough to $f$ that the corresponding
averaged pair is close. Consequently
\begin{equation}\label{eq:core}
 \overline{\{(J_hp,AJ_hp):h>0,\ p\in\mathcal P\}}
       =\operatorname{graph}A.
\end{equation}
This is the precise smoothed core assertion. It does
not identify the closure of the unsmoothed polynomial graph.

\section{Laplace resolvent on the full domain}
For $\beta>0$, contraction makes the Bochner integral
\[
 R_\beta f=\int_0^\infty e^{-\beta t}P_tf\,dt
\]
absolutely convergent, with norm at most
$\|f\|_\infty/\beta$. It is positive on real nonnegative
inputs. For a finite integral $u_T=\int_0^T
e^{-\beta t}P_tf\,dt$, the shifted-orbit identity is
\[
 P_hu_T=e^{\beta h}\int_h^{T+h}e^{-\beta s}P_sf\,ds.
\]
Differentiate at $h=0$ from the right to obtain
\begin{equation}\label{eq:finitelaplace}
 u_T\in D(A),\qquad
 Au_T=\beta u_T+e^{-\beta T}P_Tf-f.
\end{equation}
The tails of both sides converge in supremum norm.
Closedness of the graph gives
\begin{equation}\label{eq:rightinverse}
 R_\beta f\in D(A),\qquad (\beta-A)R_\beta f=f.
\end{equation}
For $u\in D(A)$, the orbit derivative from
Proposition~\ref{prop:graph} yields
\[
 \frac{d}{dt}(e^{-\beta t}P_tu)
 =-e^{-\beta t}P_t(\beta u-Au).
\]
Integrate from zero to infinity; the endpoint at infinity
vanishes by contraction. Hence
$R_\beta(\beta u-Au)=u$ as well.
This proves the full-domain inverse formula and the bound
\[
 \beta\|u\|_\infty\leq\|\beta u-Au\|_\infty.
\]
No self-adjointness was used. The next calculation shows
why it would be inappropriate to insert it.

\section{An exact nonsymmetry witness}
Under $\nu_q$, the Poisson moments satisfy
$\int z^n\,d\nu_q=q^n$ for $n\geq0$.
Use $f(z)=z$ and $g(z)=z^2$, both in the finite
Fourier generator domain. Since $q$ and $r$ are real,
\begin{align*}
 \langle f,Ag\rangle_{L^2(\nu_q)}
 &=-2q+r(1-q^2)=0,\\
 \langle Af,g\rangle_{L^2(\nu_q)}
 &=-q+\tfrac r2(q^2-1)=-2q.
\end{align*}
Here the inner product is conjugate linear in its first
argument. If $q\ne0$ the pairings differ. Thus no
Hermitian realization agreeing with this core can be
inferred from the invariant Poisson law.
Invariance, positivity and entropy decay remain valid;
they do not require reversibility.

\section{Forward densities and a stationary Bregman identity}
For an arbitrary probability measure $\mu$, Fubini gives
the positive-time density relative to $\nu_q$:
\begin{equation}\label{eq:density}
 f_t(w)=\int_\T
       \frac{p_{\phi_t(z)}(w)}{p_q(w)}\,\mu(dz).
\end{equation}
The ratio is jointly continuous and strictly positive
on the compact product circle for fixed $t>0$. It has
a positive minimum and finite maximum, independent of
$\mu$. On every compact time interval inside $(0,\infty)$,
the denominators stay uniformly away from zero. Ordinary
differentiation of the rational kernel, uniformly in
$z,w$, proves that $t\mapsto f_t$ is continuously
differentiable in supremum norm. The same denominator
bound gives a bounded angular derivative in $w$, hence
a Lipschitz bound in chord distance on $\T$.

In particular $\int f_t\,d\nu_q=1$ and
$E(t)=\int f_t\log f_t\,d\nu_q$ is finite and differentiable
for $t>0$. Positivity and the uniform bounds on compact
time intervals justify the usual scalar derivative,
\begin{equation}\label{eq:ordinaryentropy}
 E'(t)=\int_\T \dot f_t\log f_t\,d\nu_q,
\end{equation}
since $\int\dot f_t\,d\nu_q=0$.
This is a statement about actual forward laws and their
Radon--Nikodym density, not an assumed backward-generator
domain for $\log f_t$.

For positive $a,b$, set $B(a,b)=a\log(a/b)-a+b$.
It is nonnegative by convexity of $s\log s-s$.
If $f\nu_q$ is a probability law with positive continuous
$f$, invariance and the two marginals of
$\nu_q(dz)K_h(z,dw)$ give the exact identity
\begin{equation}\label{eq:bregmanstep}
\begin{split}
 &\int_{\T^2}B(f(z),f(w))\,\nu_q(dz)K_h(z,dw)\\
 &\qquad=\int_\T f\log f\,d\nu_q
        -\int_\T\log f\,d((f\nu_q)K_h).
\end{split}
\end{equation}
Indeed expansion of $B$ gives four terms. The first
integrates to $\int f\log f\,d\nu_q$; the mixed term
is the last term in \eqref{eq:bregmanstep}; and the two
linear terms cancel because both stationary marginals
are $\nu_q$. All integrands are bounded, so Fubini is
immediate.
For $f=f_t$, Chapman--Kolmogorov identifies the last
law as $\mu_{t+h}$. Divide by $h$ and use the
supremum-norm derivative of $f_t$ to conclude
\begin{equation}\label{eq:steplimit}
 \lim_{h\downarrow0}\frac1h
 \int B(f_t(z),f_t(w))\,\nu_q(dz)K_h(z,dw)=-E'(t).
\end{equation}

\section{The singular jump limit and its domination}
Let $v_h(z)=\phi_h(z)$, $z\in\T$.
Equation \eqref{eq:velocity} gives
\begin{equation}\label{eq:radial}
 \left.\frac{d}{dh}(1-|v_h(z)|^2)\right|_{h=0}=2.
\end{equation}
To check the independence of $q$, compute
$-2\operatorname{Re}(\bar z V(z))$.
The $-z$ term contributes two, while
$\bar z(1-z^2)=\bar z-z$ is purely imaginary.
The derivatives are uniform on the compact circle.
For $z\ne w$ it follows that
\begin{equation}\label{eq:jump}
 \frac{p_{v_h(z)}(w)}h\longrightarrow\frac2{|w-z|^2}.
\end{equation}
A pointwise limit alone does not justify the singular
integral in \eqref{eq:entropyderivative}; we give a
bound that also handles points arbitrarily close to the diagonal.

Uniform differentiability and \eqref{eq:radial} imply
that, for all sufficiently small $h>0$,
\[
 |v_h(z)-z|\leq C h,\qquad
 h\leq1-|v_h(z)|^2\leq3h,
\]
uniformly in $z$. Since $|w-v_h(z)|\geq1-|v_h(z)|\geq h/2$,
the triangle inequality gives
\[
 |w-z|\leq|w-v_h(z)|+Ch
       \leq(1+2C)|w-v_h(z)|.
\]
Consequently
\begin{equation}\label{eq:domination}
 h^{-1}p_{v_h(z)}(w)|w-z|^2\leq3(1+2C)^2.
\end{equation}
Suppose $f\geq\delta>0$ and
$|f(z)-f(w)|\leq L|w-z|$. Taylor's integral remainder
for $s\log s-s$ yields
\[
 0\leq B(f(z),f(w))\leq\frac{L^2}{2\delta}|w-z|^2.
\]
Combine this bound with \eqref{eq:domination} and the
positive lower bound of $p_q$ on the circle.
The functions
\[
 \frac{p_{v_h(z)}(w)}{h\,p_q(w)}B(f(z),f(w))
\]
are uniformly bounded on $\T^2$, with value zero when
$z=w$. Their pointwise limit is the singular integrand
in \eqref{eq:entropyderivative}, with the stated zero
diagonal convention. Dominated convergence proves its
integrability and the equality of its integral with the
limit in \eqref{eq:steplimit}. Apply this to $f_t$ using
the positivity and Lipschitz bounds from the preceding
section. This proves the exact production formula.

The drift contributes no separate square-gradient term
to this formula. Its effect is retained in $f_t$ and
the invariant law; \eqref{eq:radial} shows that the
off-diagonal jump rate is the Cauchy rate. No diffusion
chain rule or symmetry assumption enters this argument.

\section{A nonsharp entropy inequality and decay from initial time}
For $w\in\T$,
\[
 p_q(w)\leq\frac{1+|q|}{1-|q|},\qquad |w-z|\leq2.
\]
Thus away from the diagonal
$2/(|w-z|^2p_q(w))\geq c_q$.
Since $B(f(z),f(z))=0$, the product with $B$ satisfies
the same inequality everywhere under the zero convention.
For a normalized positive Lipschitz $f$ bounded away
from zero, Fubini gives
\begin{align*}
 \int_{\T^2}B(f(z),f(w))\,d\nu_q(z)d\nu_q(w)
 &=\int_\T f\log f\,d\nu_q-\int_\T\log f\,d\nu_q\\
 &\geq\Ent_{\nu_q}(f).
\end{align*}
The last step is Jensen's inequality for the logarithm
and $\int f\,d\nu_q=1$. We have proved the modified
logarithmic-Sobolev inequality in the exact form
\begin{equation}\label{eq:mlsi}
 c_q\Ent_{\nu_q}(f)
 \leq\int_{\T^2}\frac{2B(f(z),f(w))}{|w-z|^2p_q(w)}
            \,\nu_q(dz)\nu_q(dw).
\end{equation}
This is a jump-Bregman inequality on the displayed
positive Lipschitz class, rather than a sharp quadratic
Dirichlet-form logarithmic-Sobolev assertion.

For the actual $f_t$, the production formula implies
$E'(t)\leq-c_qE(t)$ for $t>0$. Multiplying by
$e^{c_qt}$ and integrating yields
\begin{equation}\label{eq:burnin}
 E(t)\leq e^{-c_q(t-s)}E(s),\qquad0<s\leq t.
\end{equation}
To reach initial time, we also need entropy contraction
without any initial smoothness.

If $D(\mu\Vert\nu_q)<\infty$, write $\mu=f\nu_q$.
For $h>0$ put $R_h(z,w)=p_{\phi_h(z)}(w)/p_q(w)$.
Invariance says $\int R_h(z,w)\nu_q(dz)=1$ for
$\nu_q$-almost every $w$; the row integral is also one.
The output density is
$f_h(w)=\int R_h(z,w)f(z)\nu_q(dz)$.
Apply Jensen's inequality pointwise for the nonnegative
convex function $\Psi(a)=a\log a-a+1$, then integrate:
\[
 \int\Psi(f_h)\,d\nu_q
 \leq\int\!\int R_h(z,w)\Psi(f(z))\,d\nu_q(z)d\nu_q(w)
 =\int\Psi(f)\,d\nu_q.
\]
Tonelli justifies the extended nonnegative integrals.
This is $E(h)\leq D(\mu\Vert\nu_q)$. Insert it in
\eqref{eq:burnin} and let $s\downarrow0$ to obtain
\eqref{eq:entropybound} when the initial entropy is finite.
If it is infinite the same extended-value inequality
is immediate; at $t=0$ it is equality.
Positive-time finiteness has already been proved directly
from \eqref{eq:density}, even for singular initial laws.

For any fixed $s>0$, \eqref{eq:burnin} also proves
$E(t)\to0$ for every initial probability law as
$t\to\infty$. If a probability measure satisfies
$\mu K_h=\mu$ for one $h>0$, it has finite entropy by
positive-time smoothing. Equation \eqref{eq:entropybound}
then forces $E\leq e^{-c_qh}E$, so $E=0$.
Strict convexity of $\Psi$ gives $\mu=\nu_q$.
This proves the uniqueness assertion without a path-space
construction or a reversibility premise.

\section{Classical context and exact verification scope}
The Poisson kernel and harmonic-composition construction
are classical. The pole equation
$\dot a=-a+(r/2)(1-a^2)$ is also the constant-forcing,
zero-frequency, unit-noise specialization of the
Cauchy-driven phase equation in Tanaka~\cite{tanaka},
equations (9) and (17). Its relation to a known invariant
family must therefore be acknowledged; we do not present
the rational pole flow as a new general mechanism.
The generator/resolvent arguments are standard contraction
semigroup theory. The complete kernel-to-forward-entropy
exposition and its precise formal domains are the focus here.
Priority for this explicit combined formulation has not
been established.

All mapped statements use the frozen verified baseline
\begin{center}\small\ttfamily
bab688bb6140cd6316c80680095a914ea0d346db.
\end{center}
The accompanying theorem map covers the actual Markov
kernels, Feller semigroup, closed generator, smoothed
core, Laplace resolvent, forward density, stationary
Bregman identity, singular limit, modified entropy bound
and finite nonsymmetry witness. The human proof groups
the same ingredients into harmonic and Banach-space
arguments; intermediate regroupings are not claimed as
new named Lean theorems.

The results concern the displayed circle kernels and
the continuous-function generator. A trajectory law,
maximal forward-density generator domain, unsmoothed
Fourier graph-core theorem, sharp entropy constant or
diffusion curvature estimate is not asserted here.
In particular an invariant wrapped Cauchy law cannot
justify replacing the actual nonsymmetric operator by
a symmetric diffusion.

\begin{thebibliography}{9}
\bibitem{tanaka} T. Tanaka,
Low-dimensional dynamics of phase oscillators driven by
Cauchy noise, primary preprint, arXiv:2009.04273.
\url{https://arxiv.org/abs/2009.04273}.
\end{thebibliography}
\end{document}
